% ==== Appendix B: Project Management and Financial Aspects (PO11) =====
% Sources: WRITING_PLAN section 8 row B; Modal billing for cloud spend;
% the deck's Time Plan slide as the Gantt seed. Build the Gantt with
% figures/_src/make_gantt.py (matplotlib; data chart, not architecture).
\chapter{Project Management and Financial Aspects}

\section{Project Timeline}
The project ran from February to July 2026. Figure \ref{fig:gantt} shows
the schedule: an initial phase of paper collection and study overlapping a
review-paper draft, followed by architecture design and an ablation
campaign, then the production training run and its evaluation gates, with
thesis writing in the final weeks.

\begin{figure}[!htbp]
\centering
\includegraphics[width=\textwidth]{fig_b_1_gantt}
\caption{\textbf{Project timeline.} The seven work phases from February to July 2026.}
\label{fig:gantt}
\end{figure}

\section{Expenses}
The project was carried out almost entirely on rented cloud compute, which
kept the hardware cost low. Table \ref{tab:expenses} itemizes the spend.
The dominant item is the single production training run on a cloud A100;
the many ablation and evaluation runs were individually inexpensive
because they used the smaller T4 and L40S instances for short durations.

\begin{table}[!htbp]
\centering
\caption{\textbf{Project expenses.} Cloud figures are approximate and reflect the on-demand rates of the compute provider.}
\label{tab:expenses}
\small
\begin{tabular}{lc}
\toprule
Item & Approximate cost (USD)\\
\midrule
Production training run (A100, about one day) & 10 to 15\\
Ablation sweeps (T4 and L40S, many short runs) & 15 to 25\\
Evaluation runs (T4, a fraction of a dollar each) & 5\\
Stereo camera rig & 60\\
Storage and miscellaneous cloud & 10\\
\midrule
Total & 100 to 115\\
\bottomrule
\end{tabular}
\end{table}

\section{Economic Analysis}
The economic case for the approach rests on the cost of the deployed
sensor rather than on the cost of development. A passive stereo rig
requires two commodity image sensors and an embedded accelerator, a bill
of materials on the order of a few tens of dollars, whereas a comparable
active-depth or LiDAR unit costs several times as much and adds an emitter.
Because the model runs on the device, no recurring cloud-inference cost is
incurred in deployment. The combination of a low sensor cost and
on-device inference makes the approach attractive for price-sensitive
robotics and embedded products.
